% =======================================================================
% Clase 1 — Arquitectura IA Generativa y el Ecosistema LLM
% Curso de IA Generativa · EscuelaIT · Abril 2026
% Compilación: make pdf   (xelatex)
% =======================================================================

\documentclass[aspectratio=169,11pt]{beamer}
\usetheme[progressbar=frametitle, sectionpage=progressbar, numbering=fraction]{metropolis}

% --------- Idioma y fuentes (xelatex) ---------
\usepackage{polyglossia}
\setmainlanguage{spanish}
\usepackage{fontspec}

% --------- Paquetes ---------
\usepackage{booktabs}
\usepackage{tabularx}
\usepackage{array}
\usepackage{xcolor}
\usepackage{tikz}
\usetikzlibrary{shapes.geometric, arrows.meta, positioning, fit, backgrounds}
\usepackage{listings}
\usepackage{hyperref}
\usepackage{fontawesome5}

% --------- Colores ---------
\definecolor{mBlue}{HTML}{1976D2}
\definecolor{mRed}{HTML}{D32F2F}
\definecolor{mGreen}{HTML}{388E3C}
\definecolor{mOrange}{HTML}{F57C00}
\definecolor{mGray}{HTML}{616161}
\definecolor{mLightGray}{HTML}{EEEEEE}
\definecolor{mCodeBg}{HTML}{FAFAFA}
\definecolor{usaColor}{HTML}{1976D2}
\definecolor{chinaColor}{HTML}{D32F2F}

% --------- Listings / código ---------
\lstdefinestyle{pystyle}{
  backgroundcolor=\color{mCodeBg},
  basicstyle=\ttfamily\scriptsize,
  keywordstyle=\color{mBlue}\bfseries,
  stringstyle=\color{mGreen},
  commentstyle=\color{mGray}\itshape,
  numbers=none,
  showstringspaces=false,
  frame=leftline,
  rulecolor=\color{mBlue},
  breaklines=true,
  language=Python,
  columns=fullflexible,
  keepspaces=true
}
\lstdefinestyle{bashstyle}{
  backgroundcolor=\color{mCodeBg},
  basicstyle=\ttfamily\scriptsize,
  keywordstyle=\color{mBlue}\bfseries,
  commentstyle=\color{mGray}\itshape,
  numbers=none,
  showstringspaces=false,
  frame=leftline,
  rulecolor=\color{mOrange},
  breaklines=true,
  language=bash,
  columns=fullflexible,
  keepspaces=true
}
\lstset{style=pystyle}

% --------- Metadata ---------
\title{Arquitectura IA Generativa y el Ecosistema LLM}
\subtitle{Clase 1 de 5 · Curso de IA Generativa}
\author{Jean Pierre Mandujano}
\institute{EscuelaIT}
\date{Abril 2026}

% --------- Helpers ---------
\newcommand{\hl}[1]{\textcolor{mBlue}{\textbf{#1}}}
\newcommand{\flag}[1]{\textcolor{#1}{\faFlag}}

% =======================================================================
\begin{document}

% ----------------------------------------- Portada
\maketitle

% ----------------------------------------- Bloque 0 · Apertura
\section{Apertura}

\begin{frame}{Antes de empezar — una demo}
  \begin{center}
    \vspace{0.5cm}
    {\Large \textbf{peru-elecciones.deepskill.space}}\\[0.6cm]
    {\large Asistente IA sobre datos oficiales de la ONPE}\\[1cm]
    \begin{tikzpicture}
      \node[draw=mBlue, thick, rounded corners, minimum width=10cm, minimum height=2.5cm, fill=mLightGray] (box) {
        \begin{minipage}{9.5cm}\centering
          \textit{"¿Quién va segundo en Lima?"}\\
          \textit{"¿Cuánta participación hay en el norte?"}\\
          \textit{"Dame los partidos con más crecimiento en las últimas 2 horas"}
        \end{minipage}
      };
    \end{tikzpicture}\\[0.6cm]
    \small \textcolor{mGray}{Al terminar este curso, sabrán construir esto — y mejor.}
  \end{center}
\end{frame}

\begin{frame}{Nosotros ya somos usuarios de software con IA}
  \small
  \begin{tikzpicture}[node distance=0.15cm, every node/.style={font=\small}]
    \node[draw=mGray, rounded corners, fill=mLightGray, minimum width=2.2cm, minimum height=1.3cm, align=center] (n1) {\textbf{2022}\\ChatGPT\\(navegador)};
    \node[draw=mGray, rounded corners, fill=mLightGray, minimum width=2.2cm, minimum height=1.3cm, align=center, right=of n1] (n2) {\textbf{2023}\\Copilot\\in-line};
    \node[draw=mGray, rounded corners, fill=mLightGray, minimum width=2.2cm, minimum height=1.3cm, align=center, right=of n2] (n3) {\textbf{2024}\\Cursor /\\Windsurf};
    \node[draw=mBlue, thick, rounded corners, fill=mLightGray, minimum width=2.2cm, minimum height=1.3cm, align=center, right=of n3] (n4) {\textbf{2025}\\Claude Code\\Codex · Qwen};
    \node[draw=mBlue, thick, rounded corners, fill=mBlue!10, minimum width=2.2cm, minimum height=1.3cm, align=center, right=of n4] (n5) {\textbf{2026}\\IDE orquesta\\agentes · ACP};

    \draw[->, thick, mGray] (n1) -- (n2);
    \draw[->, thick, mGray] (n2) -- (n3);
    \draw[->, thick, mGray] (n3) -- (n4);
    \draw[->, thick, mBlue] (n4) -- (n5);
  \end{tikzpicture}

  \vspace{0.7cm}
  \begin{center}
    \begin{minipage}{0.85\textwidth}\centering
      \small El mismo tipo de software que ya usamos como desarrolladores\\
      es lo que vamos a construir \textbf{para nuestros usuarios}.
    \end{minipage}
  \end{center}
\end{frame}

\begin{frame}{Marco mental del curso — tres capas}
  \centering
  \begin{tikzpicture}[
    every node/.style={font=\small, align=center},
    layer/.style={draw, rounded corners, minimum width=10cm, minimum height=1.4cm, thick}
  ]
    \node[layer, fill=mBlue!15, draw=mBlue] (l1) {\textbf{Software que administra}\\
      \footnotesize tu app · Cursor · Claude Code · peru-elecciones};
    \node[layer, fill=mOrange!15, draw=mOrange, below=0.4cm of l1] (l2) {\textbf{Capa agéntica}\\
      \footnotesize loops · tools · memoria · MCP · orquestación · guardrails};
    \node[layer, fill=mGreen!15, draw=mGreen, below=0.4cm of l2] (l3) {\textbf{LLM base}\\
      \footnotesize Claude · GPT · Gemini · DeepSeek · Qwen · GLM · Kimi\dots};

    \draw[->, thick] (l1) -- (l2);
    \draw[->, thick] (l2) -- (l3);
  \end{tikzpicture}

  \vspace{0.4cm}
  \small \textcolor{mGray}{Columna vertebral conceptual del curso entero. Volvemos varias veces.}
\end{frame}

\begin{frame}{El curso — 5 clases de 90 min}
  \small
  \begin{center}
  \renewcommand{\arraystretch}{1.25}
  \begin{tabular}{@{}rl@{}}
    \toprule
    \textbf{Clase 1} & Arquitectura IA Generativa y el Ecosistema LLM \textit{(hoy)} \\
    \textbf{Clase 2} & Prompt Engineering como sistema \\
    \textbf{Clase 3} & Patrones de arquitectura con LLMs (RAG, orquestación, memoria) \\
    \textbf{Clase 4} & Construcción de aplicaciones LLM \\
    \textbf{Clase 5} & Despliegue y producción \\
    \bottomrule
  \end{tabular}
  \end{center}

  \vspace{0.7cm}
  \begin{center}
    \begin{minipage}{0.8\textwidth}\centering
      \small Hoy: \textbf{el mapa del territorio}. Las 4 clases siguientes: construir.
    \end{minipage}
  \end{center}
\end{frame}

% ----------------------------------------- Bloque 1 · Fundamentos
\section{Fundamentos}

\begin{frame}{¿Qué es un LLM?}
  \small
  Red neuronal profunda (arquitectura \hl{Transformer}) entrenada sobre billones de tokens con un objetivo autosupervisado: \textbf{predecir el siguiente token}.

  \vspace{0.4cm}
  \begin{center}
  \renewcommand{\arraystretch}{1.2}
  \begin{tabular}{@{}p{4.2cm} p{4.5cm} p{4.5cm}@{}}
    \toprule
     & \textbf{ML clásico} & \textbf{LLM / fundacional} \\
    \midrule
    Modelo por tarea & uno por tarea & uno para cientos \\
    Cómo se adapta & fine-tuning específico & prompting en lenguaje natural \\
    Datos & etiquetados + features & texto crudo \\
    Al escalar & se satura & sigue leyes de potencia \\
    \bottomrule
  \end{tabular}
  \end{center}

  \vspace{0.3cm}
  \begin{center}
    \textcolor{mBlue}{\textit{Ya no entrenamos modelos. Los consumimos como infraestructura.}}
  \end{center}
\end{frame}

\begin{frame}{El Transformer en una idea}
  \small
  \begin{columns}[T]
    \begin{column}{0.55\textwidth}
      \textbf{Tres conceptos:}
      \begin{enumerate}
        \item \hl{Tokenización} — el texto se parte en subpalabras.
        \item \hl{Self-attention} — cada token mira a todos los demás en paralelo.
        \item \hl{Autoregresivo} — predice un token, lo concatena, repite.
      \end{enumerate}
      \vspace{0.3cm}
      \textit{"Una mesa redonda donde cada palabra escucha a todas las demás antes de hablar."}
    \end{column}
    \begin{column}{0.45\textwidth}
      \centering
      \begin{tikzpicture}[scale=0.85, every node/.style={font=\scriptsize}]
        \foreach \i/\w in {0/Hola,1/{¿cómo},2/estás,3/hoy}
          \node[draw=mBlue, rounded corners, minimum width=1.1cm, minimum height=0.5cm, fill=mBlue!10] (t\i) at (\i*1.3, 0) {\w};
        \foreach \i/\j in {0/1,0/2,0/3,1/2,1/3,2/3}
          \draw[-, mBlue!40, thin] (t\i) to [bend left=35] (t\j);
        \node[above=0.9cm of t1.north east, anchor=south, font=\scriptsize\itshape, text=mGray] {self-attention};
      \end{tikzpicture}
      \vspace{0.3cm}
      \textcolor{mGray}{\scriptsize Vaswani et al., \textit{Attention Is All You Need} (2017)}
    \end{column}
  \end{columns}
\end{frame}

\begin{frame}{Chat vs razonador — la distinción de 2025-2026}
  \small
  \begin{center}
  \renewcommand{\arraystretch}{1.2}
  \begin{tabular}{@{}p{3.3cm} p{5cm} p{5cm}@{}}
    \toprule
     & \textbf{Modelo chat} & \textbf{Modelo razonador} \\
    \midrule
    Ejemplos & GPT-4o, Sonnet, Flash, V3 & o-series, Opus w/ thinking, R1, GLM-5.1 \\
    Cómo responde & un forward pass & cadena de razonamiento larga antes \\
    Latencia a 1er token & 200--500 ms & 30--180 segundos \\
    Coste & bajo & alto (los tokens de thinking se facturan) \\
    Bueno en & chat, extracción, RAG, clasificación & matemáticas, coding agéntico, planning \\
    \bottomrule
  \end{tabular}
  \end{center}

  \vspace{0.4cm}
  \begin{center}
    \begin{minipage}{0.85\textwidth}\centering
      \small Los razonadores escalan en \textbf{tiempo de inferencia}. En 2026 casi todos los\\frontier ofrecen variantes razonadoras.
    \end{minipage}
  \end{center}
\end{frame}

\begin{frame}{Tres límites duros que no se pueden ignorar}
  \small
  \begin{center}
  \renewcommand{\arraystretch}{1.35}
  \begin{tabular}{@{}p{3.8cm} p{5cm} p{4.4cm}@{}}
    \toprule
    \textbf{Límite} & \textbf{Qué pasa} & \textbf{Mitigación} \\
    \midrule
    \hl{Alucinación} & genera plausibilidad, no verdad & RAG, tool use, verificación \\
    \hl{Conocimiento congelado} & no sabe lo posterior a su cutoff & búsqueda web, RAG fresco \\
    \hl{Cálculo / consistencia} & falla en aritmética y contextos largos & code interpreter, memoria explícita \\
    \bottomrule
  \end{tabular}
  \end{center}

  \vspace{0.6cm}
  \begin{center}
    \begin{minipage}{0.85\textwidth}\centering
      \textcolor{mBlue}{\textit{Estos límites \textbf{no son bugs a corregir}.\\Son la razón de que exista el resto del curso.}}
    \end{minipage}
  \end{center}
\end{frame}

\begin{frame}{Recordatorio — tres capas}
  \centering
  \begin{tikzpicture}[
    every node/.style={font=\small, align=center},
    layer/.style={draw, rounded corners, minimum width=9.5cm, minimum height=1.2cm, thick}
  ]
    \node[layer, fill=mBlue!15, draw=mBlue] (l1) {\textbf{Software que administra}};
    \node[layer, fill=mOrange!15, draw=mOrange, below=0.3cm of l1] (l2) {\textbf{Capa agéntica} \footnotesize (loops, tools, memoria, MCP)};
    \node[layer, fill=mGreen!15, draw=mGreen, below=0.3cm of l2] (l3) {\textbf{LLM base}};
  \end{tikzpicture}

  \vspace{0.6cm}
  \small El curso enseña a construir \textbf{capa 2} y \textbf{capa 1}, consumiendo capa 3 como commodity.

  \vspace{0.3cm}
  \textcolor{mGray}{\footnotesize Explicará, por ejemplo, por qué el chat de un proveedor hace cosas que su API no hace.}
\end{frame}

% =======================================================================
% PENDIENTE DE GENERAR:
%   Bloque 2 — Tokens, contexto y límites (~5 slides)
%   Bloque 3 — Proveedores (top 9, Llama vs chinos, chat vs API, self-host, descarga) (~10 slides)
%   Bloque 4 — Frameworks y protocolos (MCP, ACP, OAuth) (~5 slides)
%   Bloque 5 — APIs chat completions + demo base_url (~6 slides)
%   Cierre — 3 take-aways + preview + Q&A (~3 slides)
% =======================================================================

\end{document}
